\textbf{1. Kiến thức cần nhớ}
\begin{itemize}
\item \textbf{Căn bậc hai số học}: Với số thực $a \ge 0$, số $\sqrt{a}$ được gọi là căn bậc hai số học của $a$:
$$\boxed{x = \sqrt{a} \iff \begin{cases} x \ge 0 \\ x^2 = a \end{cases}}$$
\item Với $a > 0$, số $a$ có đúng hai căn bậc hai là hai số đối nhau: $\sqrt{a}$ và $-\sqrt{a}$. Số $0$ có một căn bậc hai là $0$ ($\sqrt{0} = 0$). Số âm không có căn bậc hai.
\item \textbf{So sánh căn bậc hai}: Với hai số không âm $a, b$:
$$\boxed{a < b \iff \sqrt{a} < \sqrt{b}}$$
\item \textbf{Căn thức bậc hai}: Biểu thức $\sqrt{A}$ được gọi là căn thức bậc hai của $A$ ($A$ gọi là biểu thức dưới dấu căn).
$$\boxed{\sqrt{A}\text{ xác định (có nghĩa)} \iff A \ge 0}$$
\item \textbf{Hằng đẳng thức}: Với mọi biểu thức $A$, ta có:
$$\boxed{\sqrt{A^2} = |A| = \begin{cases} A & \text{khi } A \ge 0 \\ -A & \text{khi } A < 0 \end{cases}}$$
\end{itemize}

\textbf{2. Phương pháp giải các dạng}

%% DANG 9D3?7-1
\textbf{Dạng 1. Tìm căn bậc hai, căn bậc hai số học}
\begin{enumerate}
\item Căn bậc hai số học của $a \ge 0$ là $\sqrt{a} \ge 0$.
\item Các căn bậc hai của số dương $a$ gồm hai giá trị: $\pm \sqrt{a}$.
\item Chú ý: Kí hiệu $\sqrt{a}$ luôn chỉ căn bậc hai \textbf{số học} (mang giá trị không âm).
\end{enumerate}

%% DANG 9D3?7-2
\textbf{Dạng 2. Điều kiện xác định của căn thức bậc hai}
\begin{enumerate}
\item Để $\sqrt{A}$ có nghĩa, ta đặt điều kiện: $\boxed{A \ge 0}$.
\item Để $\dfrac{1}{\sqrt{A}}$ có nghĩa, ta đặt điều kiện: $\boxed{A > 0}$.
\item Để $\dfrac{\sqrt{A}}{B}$ có nghĩa, ta đặt điều kiện: $\boxed{\begin{cases} A \ge 0 \\ B \ne 0 \end{cases}}$.
\item Giải bất phương trình tìm tập hợp giá trị của $x$ và kết luận.
\end{enumerate}

%% DANG 9D3?7-3
\textbf{Dạng 3. Hằng đẳng thức $\sqrt{A^2} = |A|$}
\begin{enumerate}
\item Đưa biểu thức dưới dấu căn về dạng bình phương một tổng hoặc một hiệu: $A^2 = (a \pm b)^2$.
\item Áp dụng hằng đẳng thức: $\sqrt{(a \pm b)^2} = |a \pm b|$.
\item Xét dấu của $a \pm b$ để bỏ dấu giá trị tuyệt đối:
\begin{itemize}
\item Nếu $a \ge b \implies |a - b| = a - b$.
\item Nếu $a < b \implies |a - b| = -(a - b) = b - a$.
\end{itemize}
\end{enumerate}

%% DANG 9D3?7-4
\textbf{Dạng 4. Tính giá trị, so sánh căn bậc hai}
\begin{enumerate}
\item Để so sánh hai số (một số có căn, một số là số hữu tỉ): ta bình phương hai số hoặc đưa số hữu tỉ vào trong dấu căn ($m = \sqrt{m^2}$ với $m \ge 0$) rồi so sánh các biểu thức dưới dấu căn.
\item Để so sánh tổng hai căn với một số, ta có thể so sánh bình phương của hai vế.
\end{enumerate}
